\editionPage[\pdfbookmark[0]{Indications des défis}{bm-indices-debut-parent}]{indices-debut}{Indications progressives}{Défis 1 à 3}{Ne lire que l’aide nécessaire.}
\begin{blocEdition}{indices-debut}
\editionEtape{D01 · Cent termes, sans cent calculs}\textbf{Indice 1.} Regrouper les termes par paires consécutives.\par\medskip\textbf{Indice 2.} Chaque paire $-(2k-1)+2k$ vaut 1.\par
\editionEtape{D02 · Retrouver tout le budget}\textbf{Indice 1.} Exprimer chaque reste comme fraction du budget initial.\par\medskip\textbf{Indice 2.} Le reste final représente $\frac23\times\frac34\times\frac25$ du budget.\par
\editionEtape{D03 · Jusqu’où peut-on doubler ?}\textbf{Indice 1.} Après $n$ étapes, combien de facteurs 2 a-t-on introduits ?\par\medskip\textbf{Indice 2.} Calculer $2^{10}$, puis comparer $2^{19}$ et $2^{20}$ à un million.\par

\end{blocEdition}
\editionPage{indices-4}{Indications progressives}{Défis 4 à 6}{Ne lire que l’aide nécessaire.}
\begin{blocEdition}{indices-4}
\editionEtape{D04 · Racine de deux, un nombre non décimal}\textbf{Indice 1.} Encadrer $\sqrt2$ entre deux entiers consécutifs grâce à $1^2$ et $2^2$.\par\medskip\textbf{Indice 2.} Écrire les carrés de 1 à 9 et regarder leur chiffre des unités.\par
\editionEtape{D05 · Comparer sans choisir de valeur}\textbf{Indice 1.} Développer les deux expressions d’aire.\par\medskip\textbf{Indice 2.} Calculer l’aire de A moins celle de B.\par
\editionEtape{D06 · Compter les pièces}\textbf{Indice 1.} Appeler $x$ le nombre de pièces de 2 euros.\par\medskip\textbf{Indice 2.} Il reste alors $42-x$ pièces de 1 euro.\par

\end{blocEdition}
\editionPage{indices-7}{Indications progressives}{Défis 7 à 9}{Ne lire que l’aide nécessaire.}
\begin{blocEdition}{indices-7}
\editionEtape{D07 · Partager un bénéfice}\textbf{Indice 1.} Diviser les trois sommes investies par 6 000.\par\medskip\textbf{Indice 2.} La première année, le bénéfice se partage en $2+3+5=10$ parts.\par
\editionEtape{D08 · Annuler une remise}\textbf{Indice 1.} Une baisse de 20\,\% correspond au coefficient multiplicateur 0,8.\par\medskip\textbf{Indice 2.} Chercher le coefficient $k$ tel que $0{,}8\times k=1$.\par
\editionEtape{D09 · Un découpage variable}\textbf{Indice 1.} L’aire retirée est $4x^2$.\par\medskip\textbf{Indice 2.} Si l’aire restante est 116, chaque carré retiré a une aire de 16.\par

\end{blocEdition}
\editionPage{indices-10}{Indications progressives}{Défis 10 à 12}{Ne lire que l’aide nécessaire.}
\begin{blocEdition}{indices-10}
\editionEtape{D10 · Même moyenne, autre médiane}\textbf{Indice 1.} Les deux sommes valent 60.\par\medskip\textbf{Indice 2.} Avec six valeurs, observer les valeurs de rangs 3 et 4.\par
\editionEtape{D11 · Un jeu de dés équitable ?}\textbf{Indice 1.} Construire le tableau des 36 issues et marquer celles qui font gagner Léa.\par\medskip\textbf{Indice 2.} Un produit est impair seulement si les deux nombres sont impairs : utiliser l’événement contraire.\par
\editionEtape{D12 · Six triangles de même aire}\textbf{Indice 1.} Noter $a$, $b$ et $c$ les aires de GBI, GCJ et GAK : $[GI]$, $[GJ]$ et $[GK]$ sont des médianes de petits triangles.\par\medskip\textbf{Indice 2.} La médiane $[AI]$ partage ABC en deux parties de même aire.\par

\end{blocEdition}
\editionPage{indices-13}{Indications progressives}{Défis 13 à 15}{Ne lire que l’aide nécessaire.}
\begin{blocEdition}{indices-13}
\editionEtape{D13 · Le pont le plus court}\textbf{Indice 1.} Faire glisser B vers la rivière, perpendiculairement aux rives, d’une longueur égale à la largeur de la rivière.\par\medskip\textbf{Indice 2.} Dans le parallélogramme MNBB’, comparer NB et MB’, puis chercher le plus court chemin de A à B’.\par
\editionEtape{D14 · Varignon et les aires}\textbf{Indice 1.} Dans le triangle ABD, comparer les aires de AID et ABD, puis de AIL et AID.\par\medskip\textbf{Indice 2.} Les quatre triangles des coins représentent ensemble la moitié de l’aire de ABCD.\par
\editionEtape{D15 · Le câble et ses attaches}\textbf{Indice 1.} Calculer d’abord la distance entre les points d’attache.\par\medskip\textbf{Indice 2.} Deux extrémités demandent deux réserves de 0,75 m.\par

\end{blocEdition}
\editionPage{indices-16}{Indications progressives}{Défis 16 à 18}{Ne lire que l’aide nécessaire.}
\begin{blocEdition}{indices-16}
\editionEtape{D16 · Le plus grand triangle du demi-cercle}\textbf{Indice 1.} Le triangle AMB est rectangle en M.\par\medskip\textbf{Indice 2.} La hauteur issue de M est au plus égale au rayon OM.\par
\editionEtape{D17 · Une pyramide de verre}\textbf{Indice 1.} Le volume d’une pyramide est le tiers du produit de l’aire de la base par la hauteur.\par\medskip\textbf{Indice 2.} Le triangle SOM est rectangle en O et $OM$ vaut la moitié du côté de la base.\par
\editionEtape{D18 · La suite de Syracuse}\textbf{Indice 1.} Écrire n après chaque tour dans un tableau.\par\medskip\textbf{Indice 2.} Après 1, on obtient 4, puis 2, puis 1.\par

\end{blocEdition}
